\section{The bound pair and the order of events}
\label{sec:pair}

\subsection{The moment of binding}

Two fermions in the Yukawa potential (\ref{eq:yukawa}) form a bound state if and only if the number
$\DVS$ of Eq.~(\ref{eq:DC}) exceeds a critical value, $D_c=1.67981$. On a deflating spacetime
$\DVS\propto e^{\td}$, so that the pair becomes bound at a definite time. Measure the clock from that
moment:
\begin{equation}
\tpair=\ln\frac{\DVS}{D_c}=\gsum\,\big(x4-x4_{\mathrm{binding}}\big).
\label{eq:tpair}
\end{equation}
Table~\ref{tab:dimer} and Fig.~\ref{fig:binding} show the scattering length $a_{\mathrm{sc}}$ of two
fermions and the binding energy $E_b$ of the pair as the clock passes the binding.

\begin{table}
\caption{\label{tab:dimer}Two fermions as the clock passes the moment of binding: the scattering
length in units of $1/m_s$, the number of bound states, and the binding energy in units of
$m_s^{2}/m$.}
\begin{ruledtabular}
\begin{tabular}{ccccc}
$\tpair$ & $\DVS$ & $a_{\mathrm{sc}}\,m_s$ & bound states & $E_b\,m/m_s^{2}$ \\
\hline
{{ROWS:tbl_dimer}}
\end{tabular}
\end{ruledtabular}
\end{table}

\emph{The scattering length.} Just before the binding the scattering length is large and negative,
just after it large and positive: $-134.0/m_s$ at $0.99\,D_c$ and $+135.9/m_s$ at $1.01\,D_c$. At the
moment of binding it is infinite: the two fermions are at a scattering resonance. \emph{Every pair
of fermions of the model passes through this resonance once, at the same time.}
Section~\ref{sec:cosmo} examines what that means if the fermions are the dark matter.

\emph{The binding energy.} Just after the binding the binding energy starts quadratically,
\begin{equation}
E_b\ \propto\ \big(e^{\tpair}-1\big)^{2}\,\frac{m_s^{2}}{m},
\label{eq:Eb}
\end{equation}
with the fitted exponent 2.01 between 2 and 8 per cent above $D_c$. Because $E_b$ vanishes at the
threshold, the time $1/E_b$ that the pair needs to adjust is infinite there. However slow the
deflation is, the pair cannot follow it adiabatically at first.

\emph{The same threshold from the gap equation.} In vacuum the linearized gap equation is the
equation of a bound state of two fermions, with a relativistic kernel. It reaches the eigenvalue 1 at
$\DVS=1.7347$, $1.6902$ and $1.6810$ for $m_s/m=0.3$, $0.1$ and $0.03$, approaching the value
$D_c$ of the Schr\"odinger equation as the radial mode becomes light and the binding
non-relativistic.

\begin{figure*}
\includegraphics[width=\textwidth]{fig_binding}
\caption{\label{fig:binding}The bound pair along the clock. (a) Scattering length of two fermions in
units of $1/m_s$ against $\tpair$: the resonance at the moment of binding. (b) Binding energy of the
pair. (c) Cross-section per mass of two fermions at three relative velocities while the clock passes
the binding, for the example particle of Sec.~\ref{sec:dm}.}
\end{figure*}

\subsection{The order of events}

The other thresholds of the model are thresholds in the number $C$ of Eq.~(\ref{eq:DC}),
which grows at the same rate as $\DVS$. Because $C/\DVS=4\pi m/m_s$ is fixed, the order in which the
events happen depends on the ratio $m_s/m$ alone. Table~\ref{tab:events} gives the clock reading
$\tpair$ of each event; a negative number means that the event precedes the binding of the pair. The
events other than the binding are derived in Sec.~\ref{sec:cjt}.

\begin{table*}
\caption{\label{tab:events}The order of events along the deflation: the clock reading $\tpair$ of
each event, relative to the moment at which two fermions form a bound pair, for six ratios of the
mass $m_s$ of the radial mode to the fermion mass $m$.}
\begin{ruledtabular}
\begin{tabular}{p{0.42\textwidth}lcccccc}
event & threshold & $0.1$ & $0.3$ & $1$ & $3$ & $10$ & $30$ \\
\hline
{{ROWS:tbl_events}}
\end{tabular}
\end{ruledtabular}
\end{table*}

For a \emph{light radial mode}, $m_s/m\le3$ in the table, the pair becomes bound only \emph{after}
the one-loop potential of the Vorton-Scalar has fallen below its vacuum value at a field of three
times the vacuum value: the deflation endangers the vacuum of the model before it binds the pair. The
condition for the opposite order is $4\pi D_c\,m/m_s<4.27$, that is, $m_s/m>4.9$. For a \emph{heavy
radial mode}, $m_s/m=10$ or $30$, the pair becomes bound first, and matter becomes self-bound and
loses its fermion mass later. Nothing in the model decides between the two orders; the table states
what follows from each.

